\begin{frontmatter}

\title{Orthogonal Defect Classification for \textit{Defects4J} using an\\
       LLM-Driven Scientific Approach}

% \author[iut]{Hasin Mahtab Alvee\corref{cor1}}
% \ead{hasin.alvee@iut-dhaka.edu}

% \author[iut]{Mohammad Nahiyan Kabir}
% \ead{nahiyan.swe@iut-dhaka.edu}

% \author[iut]{Md.\ Sakib Hossain}
% \ead{sakib.hossain@iut-dhaka.edu}

% \cortext[cor1]{Corresponding author.}

\affiliation[iut]{%
  organization={Department of Computer Science and Engineering,
                Islamic University of Technology},
  addressline={Board Bazar, Gazipur},
  city={Dhaka},
  postcode={1704},
  country={Bangladesh}}

%% ─── Abstract ────────────────────────────────────────────────────────────────
\begin{abstract}
Defect classification turns isolated bug reports into process-level signals that
guide debugging strategy, resource allocation, and quality improvement.
\textit{Defects4J} is the standard empirical benchmark for Java fault studies,
yet no prior work has labelled its full active set with a principled,
semantic taxonomy: existing efforts cover only tens to a few hundred bugs,
rely on structural or patch-shape taxonomies that describe \emph{how} a bug
manifests rather than assigning it a semantic defect type, or depend on
post-fix code, which is unavailable during real triage.

We close that gap. We build an automated pipeline that classifies
\textit{Defects4J} bugs into
Orthogonal Defect Classification (ODC) Design/Code defect types using Large
Language Model (LLM) reasoning. Instead of relying on a single prompt, we
structure classification as a two-variable experimental design: a
\emph{taxonomy} axis (the model's own words; the seven ODC types as a forced
choice; or the seven types plus an explicit ``Other'' escape category) crossed
with a \emph{strategy} axis (an unstructured zero-shot baseline; a strong
few-shot static prompt; or an enforced, iterative reasoning loop inspired by
scientific debugging that commits to a hypothesis, then probes held-back
evidence to confirm or refute it). The open escape category lets us measure
taxonomy coverage instead of assuming it, and the enforced loop grounds every
decision in evidence the model must actively request. We run classification in
two evidence modes, pre-fix (symptoms only) and post-fix (with the
buggy-to-fixed diff), which lets us evaluate self-consistency without needing
ground-truth labels.

To find out whether it works, we formalise five research questions spanning
defect-type distribution, taxonomy coverage, four-level classification accuracy,
the contribution of each pipeline component, and pre-fix/post-fix divergence,
and we operationalise each through a reproducible study command over a
project-balanced sample of the benchmark. A development-validation run confirms
that our pipeline and its analysis layer produce correct, well-formed metrics
end to end; a large-scale confirmatory run populates the final quantitative
results.
\end{abstract}

\begin{keyword}
Defects4J \sep
Orthogonal Defect Classification \sep
Large Language Models \sep
Scientific Reasoning \sep
Bug Taxonomy \sep
Prompt Engineering \sep
Empirical Software Engineering
\end{keyword}

\end{frontmatter}
